% 第三部分 A：摘要（由 main_zh.tex 引入）。
用策略能看到的测试来奖励它，它可能学会满足这些测试，而不是解决问题。我们在多轮代码修复中测量这一点：Qwen3.5-4B-Base 写一个函数，看到第一个失败的示例测试后重写，最多三次；由留出测试判定最终程序是否正确。离线时，基座策略通过可见测试的尝试中有 \NUM{cs_hack_all} 在留出测试上失败；\kev{} 和一个同样大小的自回归评判模型都把这类尝试排在正确尝试之下（AUROC 分别为 \NUM{csauc_all_kevg} 和 \NUM{csauc_all_arg}），并且几乎完美地把对测试的硬编码查表程序与参考解答区分开。随后我们用 GRPO 训练 100 步，每种奖励两个种子：只用可见测试；可见测试乘以 \kev{} 对代码是否 genuine 的判断；以及作为预言机的全部测试。我们没有看到可见测试奖励被利用的迹象：解出题目的比例在 KodCode val（从 \NUM{cgszmean_kod_allpass}\% 升至 \NUM{cgall_vis_kod}\%）、MBPP+ 和 HumanEval+ 上都提高了，通过可见测试的程序中错误者的比例在三个集合中的两个上降低了约一半，抄写可见测试中数值的程序仍和基座策略中一样少见。评判因子带来的变化没有超出种子之间的差异；从不奖励错误程序的预言机奖励，让错误的通过略有减少，但没有解出更多题目。本版取代了结果并非由实验产生的早期版本。代码、日志和本文的每一个数字都在仓库中。
